\chapter{常用谱对象。}
\VerdierPhysicalPageAnchor{0224}%
\VerdierSemanticAnchor{ch4d-4}
\begingroup
\small
\setlength{\parskip}{0pt}
\setlength{\emergencystretch}{1.5em}
\setlength{\abovedisplayskip}{4pt plus 2pt minus 2pt}%
\setlength{\belowdisplayskip}{4pt plus 2pt minus 2pt}%

\VerdierSemanticAnchor{ch4d-4-1}
\section{第一典范谱对象。}

\VerdierSemanticAnchor{ch4d-4-1-1}
\noindent\textbf{4.1.1.}\quad
回顾，$\widetilde{\mathbb Z}$ 表示有序集
$\mathbb Z\cup\{-\infty,+\infty\}$，也表示与这个有序集相关联的范畴
（第二章，4.3.1）。设 $\mathcal A$ 是加法范畴，并设
$Y\colon\widetilde{\mathbb Z}\longrightarrow
\operatorname{comp}(\mathcal A)$ 是函子：

\VerdierSemanticAnchor{ch4d-4-1-1-1}
\begin{equation}
\begin{aligned}
q&\in\widetilde{\mathbb Z},
  &q&\longmapsto Y(q),\\
(p,q)&\in\mathcal{F}\ell_1(\widetilde{\mathbb Z}),
  &u(p,q)&\colon Y(p)\longrightarrow Y(q).
\end{aligned}
\tag{4.1.1.1}
\end{equation}

在第一章 (3.3.6) 中，已经把一个取值于 $K(\mathcal A)$ 的谱对象
（其类型为 $\widetilde{\mathbb Z}$）（第二章，4.1.2 与 4.1.3）关联给函子
$Y$；现在回顾其构造。

首先，对 $\mathcal{F}\ell_1(\widetilde{\mathbb Z})$ 的每个元素 $(p,q)$
（即对满足 $p\leq q$ 的每一对
$p,q\in\widetilde{\mathbb Z}$），令 $\mathcal Y(p,q)$ 为态射
$u(p,q)$ 的映射锥（第一章，3.1.2）：

\VerdierSemanticAnchor{ch4d-4-1-1-2}
\begin{equation}
\begin{aligned}
\mathcal Y(p,q)^i&=Y(q)^i\oplus Y(p)^{i+1},\\
d_{\mathcal Y(p,q)}^i
&=
\begin{pmatrix}
d_{Y(q)}^i & u(p,q)^{i+1}\\
0 & -d_{Y(p)}^{i+1}
\end{pmatrix}.
\end{aligned}
\tag{4.1.1.2}
\end{equation}

对 $\mathcal{F}\ell_1(\widetilde{\mathbb Z})$ 中任意两个元素 $(p,q)$ 与
$(p',q')$，当 $p\leq p'$ 且 $q\leq q'$ 时，定义

\VerdierSemanticAnchor{ch4d-4-1-1-3}
\begin{equation}
\Phi\bigl((p,q),(p',q')\bigr)\colon
\mathcal Y(p,q)\longrightarrow\mathcal Y(p',q'),
\tag{4.1.1.3}
\end{equation}
为如下复形态射的同伦类：

\VerdierSemanticAnchor{ch4d-4-1-1-4}
\begin{equation}
\Psi\bigl((p,q),(p',q')\bigr)\colon
\mathcal Y(p,q)\longrightarrow\mathcal Y(p',q'),
\tag{4.1.1.4}
\end{equation}
该态射由下式定义：

\VerdierSemanticAnchor{ch4d-4-1-1-5}
\begin{equation}
\Psi\bigl((p,q),(p',q')\bigr)^i
=
\begin{pmatrix}
u(q,q')^i & 0\\
0 & u(p,p')^{i+1}
\end{pmatrix}.
\tag{4.1.1.5}
\end{equation}
\endgroup
